\documentclass[11pt,a4paper]{article}
\usepackage[T1]{fontenc}
\usepackage{lmodern}
\usepackage[margin=22mm,headheight=15pt]{geometry}
\usepackage{microtype,xcolor,booktabs,tabularx,array,enumitem,xurl}
\usepackage[hidelinks]{hyperref}
\usepackage{fancyhdr,titlesec}
\definecolor{studyink}{RGB}{26,49,64}
\titleformat{\section}{\large\bfseries\color{studyink}}{\thesection}{0.65em}{}
\titlespacing*{\section}{0pt}{11pt}{5pt}
\setlength{\parindent}{0pt}\setlength{\parskip}{5pt}
\setlength{\emergencystretch}{2em}
\clubpenalty=10000\widowpenalty=10000
\setlist[enumerate]{leftmargin=1.8em,itemsep=3pt,topsep=3pt}
\newcolumntype{Y}{>{\raggedright\arraybackslash}X}
\newcommand{\src}[1]{\textsuperscript{[#1]}}
\pagestyle{fancy}\fancyhf{}
\lhead{\small Student Work Assistant / Group 09}
\rhead{\small Strategic Recommendations}
\cfoot{\small\thepage}\renewcommand{\headrulewidth}{0.3pt}
\hypersetup{pdftitle={Student Work Assistant - Strategic Recommendations},pdfauthor={DSS5105 Group 09}}
\begin{document}
\begin{center}
{\LARGE\bfseries Student Work Assistant}\par
\vspace{3pt}{\large Strategic Recommendations}\par
\vspace{3pt}{\small DSS5105 Group 09 \quad | \quad Higher education and student services}
\end{center}

\section{Strategic focus and opportunity selection}
Develop three focused student services: review course changes, confirm group
commitments, and follow up multi-step administrative requirements. Each turns
unstructured messages into a source-linked decision the student can check.
GenAI interprets the language; explicit states and student review control what
becomes an obligation. These are proposed extensions of Student Work Assistant.

The target is the handoff between receiving information and knowing what to do
next. Table~\ref{tab:opportunities} ranks the opportunities by relevance to the
current user, consequence of misunderstanding, and reuse of existing components.
The initial versions use selected emails, text-based documents and pasted
notes. They require no new model training, LMS connection or campus-wide rollout.

\begin{table}[!ht]
\centering\small\renewcommand{\arraystretch}{1.12}
\begin{tabularx}{\linewidth}{@{}>{\raggedright\arraybackslash}p{0.23\linewidth}YY@{}}
\toprule
\textbf{Priority / user} & \textbf{Decision worth improving} & \textbf{Smallest useful implementation} \\
\midrule
1. Course changes / students managing several courses & Which deadline or requirement is current, and what changed? & Link messages to a deliverable; show old and proposed values beside their sources. \\
2. Group commitments / project teams & Who accepted which action, by when? & Extract actions from imported notes; record explicit acknowledgement before personal planning. \\
3. Administrative follow-up / incoming students & Which required steps remain open after a notice or receipt? & A personal checklist linking each step to an instruction and completion evidence. \\
\bottomrule
\end{tabularx}
\caption{Three student-service opportunities, ordered by fit and component reuse.}
\label{tab:opportunities}
\end{table}

\section{Opportunity 1: review course changes}
\textbf{User and demand evidence.} Students managing several courses need to
identify relevant instructions before deciding which date applies. A qualitative
University of Leeds study interviewed staff and 17 students. Students described
excessive, duplicated or irrelevant email and difficulty finding important
information; they wanted shorter, more relevant communication\src{1}. This
supports a focused course-action view rather than another general inbox summary.

\textbf{Existing alternative and specific need.} A student can copy a date into
Google Tasks, use a planner, or search the original thread. The proposed service
addresses the next question: when another notice arrives, which existing
obligation does it change? Our saved real-mail evaluation contains five duplicate
entries among 25 scored dated tasks; the sample also contains an explicit
24-hour extension\src{5}. These examples identify a concrete reconciliation
problem to solve.

\textbf{Expected impact.} Make superseded dates and unresolved contradictions
visible before the student commits study time. The useful output is one reviewed
obligation with its revision history. Assess its value through source-lookup
and correction time, missed requirements and duplicate tasks.

\textbf{Low-effort implementation.} Reuse mail intake, source quotes, task review
and the planner. Add a deliverable identifier and revision records. GenAI proposes
whether two messages concern the same work and explains the change; the student
accepts the link and authoritative date. The preview then shows the effect on
the plan. Start with one course and selected notices, including extensions and
receipts. This isolates the new work to linking, revision display and review.

\clearpage
\section{Opportunity 2: confirm group commitments}
\textbf{User and demand evidence.} Student project teams need an agreed action
list after discussions. Wang et al.'s 2025 study includes a student-led review of
group-project practice at Imperial College London. Interviews identified unclear
tasks, mismatched expectations, coordination difficulties and scheduling
conflicts. One reported case required the group to rework an assignment two days
before submission following ineffective task division\src{2}. This is evidence
of a coordination problem with consequences for deadline pressure and rework.

\textbf{Existing alternative and specific need.} Teams can maintain a shared
task board or use meeting notes. Microsoft Facilitator already captures meeting
information, and its Planner task-tracking integration is in public preview.
The proposed service focuses on the review step: distinguish a suggestion from
a decision, retain unresolved questions, and obtain acknowledgement from the
named owner before adding work to a personal plan.

\textbf{Expected impact.} Detect disagreements about responsibility before they
become last-minute rework. The output is a short list of accepted actions,
owners, dates and source quotations. Track incorrect owner assignments,
unacknowledged actions and the effort needed to prepare the next meeting's
status update. Acknowledgement documents an agreement; it does not establish
how much work a member eventually contributes.

\textbf{Low-effort implementation.} Reuse the group-text importer, source
references, review controls and planner. Add proposed, accepted, completed and
superseded states. In the first version, the team reviews an imported discussion
together and the student records each owner's acknowledgement. This avoids
building a meeting bot or a multi-user workspace. Accepted work is then offered
in the existing personal planning preview.

\textbf{Route to users.} Recruit volunteer course-project teams and approach
module coordinators to sponsor a class pilot. Students review the actions and
keep control over which accepted work enters their own plans.

\section{Opportunity 3: follow up administrative requirements}
\textbf{User and demand evidence.} Incoming students must turn instructions from
different offices into completed steps. Corwin et al.'s 2026 study draws on a
three-year research-practice partnership across two US campuses. The authors'
research summary reports that redundant or unclear email can obscure actionable
messages about registration, financial aid and academic standing\src{3}.
Locally, NUS describes onboarding as an email-invited process with separate
check-in and onboarding steps; graduate student-card instructions also come
from programme administrators\src{4}. These sources support a practical need
for following a process through, beyond reading its first notification.

\textbf{Existing alternative and specific need.} Institutional portals remain
the authority. Navigate360 already provides school to-dos, holds and
notifications at partner institutions. The proposed personal view links the
student's selected notices and receipts to individual checklist steps. A receipt
for one step must not close every requirement in the process.

\textbf{Expected impact.} Help the student see what is still pending and which
office or source to consult next. The target is fewer overlooked required steps
and fewer unnecessary repeat submissions. Measure correct step states, missed
requirements and time spent reconstructing the process from messages.

\textbf{Low-effort implementation.} Seed one checklist from a published guide
and let the student review it. Reuse extraction, evidence links and draft review;
add step IDs and links to acknowledgements. GenAI proposes which step a message
concerns. Completion comes from explicit evidence or student confirmation. The
first version stores step labels and dates, excludes identity and medical
attachments, and leaves submissions and official status changes in the portal.

\clearpage
\section{Competitive position and adoption}
Existing tools cover planning, meeting capture and campus reminders.
Table~\ref{tab:competition} locates our focus: reviewing which version applies,
who accepted work and which steps remain pending. Product names link to the
official documentation used in the comparison.

\begin{table}[!ht]
\centering\small\renewcommand{\arraystretch}{1.12}
\begin{tabularx}{\linewidth}{@{}>{\raggedright\arraybackslash}p{0.26\linewidth}YY@{}}
\toprule
\textbf{Alternative} & \textbf{Documented offer} & \textbf{Recommended focus} \\
\midrule
\href{https://support.google.com/mail/answer/9920317?hl=en}{Google Tasks} / \href{https://support.microsoft.com/en-gb/outlook/frequently-asked-questions-about-copilot-in-outlook}{Outlook Copilot} & Mail-to-task or email/calendar assistance. & Source-linked review when a later notice changes an existing obligation. \\
\href{https://mystudylife.com/}{MyStudyLife Scout} / \href{https://reclaim.ai/pricing}{Reclaim} & Weekly study planning or automatic scheduling; Reclaim offers a free tier. & Settle the obligation and its evidence before scheduling it. \\
\href{https://learn.microsoft.com/en-us/microsoftteams/facilitator-teams}{Teams Facilitator} & Collaborative meeting notes; Planner task tracking in public preview. & Review acceptance of ownership across imported discussions. \\
\href{https://play.google.com/store/apps/details?hl=en-US\&id=com.eab.se}{Navigate360 Student} & School to-dos, holds and notifications, dependent on an institutional partnership. & A student-controlled checklist linked to selected instructions and receipts. \\
\bottomrule
\end{tabularx}
\caption{Coverage described by selected product documentation, accessed 13 November 2026.}
\label{tab:competition}
\end{table}

\textbf{Positioning.} Use the project's evidence links, local storage and
content-bound approval\src{5} to make obligation changes, ownership and
completion inspectable.

\textbf{Adoption and funding.} Offer free additions to the existing assistant,
starting with selected imports and setup help at one course or onboarding
session. Approach programme offices, module coordinators and student-support
teams as pilot sponsors. Base any later support pricing on measured setup costs
and user demand.

\textbf{Implementation priority.} Preserve candidate status and date roles and
reject unsafe suggestions before adding cross-message linking; the current
handoff can promote conditional candidates to required work\src{5}. Share
evidence, review and state history across services. Use payload review, minimal
data retention and content-bound approval for external writes.

\textbf{Decision test.} After the reliability fixes, compare both versions on the
same fresh, independently labelled examples. Plan matched review tasks with
8--12 volunteers against their existing workflow, counting setup, import,
lookup, correction and confirmation time. Continue if median paired active time
falls by at least 30\% without more missed required actions; otherwise narrow
or stop the feature. This is a proposed pilot rule.

\section*{Research sources}
{\small
The research below supplies the demand evidence; linked product documentation in
Table~\ref{tab:competition} supplies the competitor evidence.
\begin{enumerate}[label={[\arabic*]},leftmargin=2.1em,itemsep=2pt]
\item Simpson, J. et al. (2025). \emph{Spamming our students? The use of email as a mass communication tool in higher education}. Staff and student interview evidence. \href{https://eprints.whiterose.ac.uk/id/eprint/221158/}{University of Leeds repository and full paper}.
\item Wang, Z. et al. (2025). \emph{Group project practices and guidance in higher education contexts}. Student-led interview findings and cases. \href{https://www.frontiersin.org/journals/education/articles/10.3389/feduc.2025.1614627/full}{Frontiers in Education}.
\item Corwin, Z. et al. (2026). \emph{Email Overload: Rethinking University Communication with At-Promise Students}. \href{https://doi.org/10.1080/19496591.2026.2639502}{Research article}; \href{https://pass.pullias.usc.edu/news/email-overload-rethinking-university-communication-with-at-promise-students/}{USC PASS author summary}.
\item NUS Registrar. \emph{Administrative Matters: Onboarding Procedures}. Official process and graduate student-card instructions. \href{https://www.nus.edu.sg/registrar/academic-activities/onboarding/administrative-matters}{NUS guidance}, accessed 13 November 2026.
\item Group 09. \emph{Student Work Assistant evaluation results}. Saved application evidence, implementation limits and version history. \href{https://github.com/YichenLi-02/student-work-assistant/blob/1de2f23/docs/evaluation/RESULTS.md}{Repository report}.
\end{enumerate}}
\end{document}
